% ==============================================================================
% EDITOR'S ROUND 3 PATCH — Exact LaTeX edits for wsc26paper.tex (V3)
% ==============================================================================
% Instructions: Apply these edits sequentially. Line numbers reference the
% current paper.tex (770 lines total).
%
% Patch format:
%   EDIT [N]: [description]
%   LOCATION: [line range]
%   ACTION: [replace|insert|append]
%   OLD: [original text]
%   NEW: [replacement text]
% ==============================================================================

================================================================================
EDIT 1: Add concrete divergence example to §1.1 (Coupling Code Problem)
================================================================================
LOCATION: After line 141 (after "...potentially inconsistent with both.")
ACTION: insert
TEXT:

In practice, silent divergence is difficult to detect. For example, consider an
experiment where the tournament scorer reads the final step's biomass while the
evolutionary algorithm reads the episode mean biomass. The evolved champion may
become the policy with volatile-but-high peaks rather than sustainable mean
production---a consistently-applied but undetected inconsistency that could
persist through publication. Such divergences are common precisely because
\emph{no single code path} connects metric definition to all three consumers.

================================================================================
EDIT 2: Add new subsection 3.5 (Design Constraints and Trade-offs)
================================================================================
LOCATION: After line 295 (after Statistical Infrastructure subsection ends)
ACTION: insert
TEXT:

\subsection{Design Constraints and Trade-offs}
\label{sec:design-constraints}

The uniform metric contract \texttt{metrics\_episode()} deliberately implements
a minimalist interface: \texttt{dict[str, float]}. This design choice reflects a
composability-versus-expressiveness trade-off. We justify this constraint and
enumerate what is intentionally excluded.

\textbf{Minimalism enables composition.} A richer contract---e.g.,
\texttt{dict[str, list[float]]} permitting per-step time series, or
\texttt{dict[str, dict[str, float]]} enabling agent-level distributions---would
require negotiation between model internals (which may naturally produce
heterogeneous outputs) and the optimization, tournament, and inference layers
(which must aggregate before comparison). By constraining the contract to scalar
floats, HEAS avoids entangling aggregation logic with metric definition. A model
that produces agent-level data can aggregate once at the \texttt{metrics\_episode()}
boundary, then reuse that scalar summary everywhere downstream. Richer contracts
would require per-stage aggregation choices, reintroducing the coupling-code
problem HEAS targets.

\textbf{What is lost.} The scalar contract does not directly represent
time-series outputs (volatility patterns, transient dynamics) or agent-level
heterogeneity (the distribution of firm sizes, not just mean firm profit). For
policy-ranking purposes---the intended use case---this loss is bounded: selecting
a regulatory regime requires aggregate summaries (mean welfare, Gini coefficient)
as the basis for comparison, even if within-regime agent-level structure remains
unmeasured. Extensions (e.g., a secondary analysis stream producing full
distributions) are possible without modifying the framework contract.

\textbf{NSGA-II as reference implementation.} Section~3.2 adopts NSGA-II as the
default multi-objective evolutionary algorithm. This choice reflects three
criteria: (1)~established performance on policy-search benchmarks in the
simulation-optimization literature; (2)~well-understood parametrization at WSC;
(3)~ease of replacement via the pluggable \texttt{run\_ea()} interface. NSGA-II's
Pareto-ranking and crowding-distance preservation provide a reasonable baseline
for two- to four-objective problems common in ABM policy design. The framework
does not depend on this choice: Section~5.4 demonstrates that algorithm
selection should be landscape-aware, and HEAS's agnostic interface enables this
flexibility without code changes.

================================================================================
EDIT 3: Add falsification criteria to §5.1 (Evaluation Questions)
================================================================================
LOCATION: After line 393 (after "...without modification.")
ACTION: insert
TEXT:

For RQ1 and RQ5 (framework verification), falsification would require evidence
that: (1)~the metric contract does not eliminate divergence across optimization,
tournament, and CI pathways, or (2)~substantial boilerplate remains after
removing model-specific logic. For RQ1, we measure lines of code directly
(coupling code only, model dynamics excluded) and compare three implementations
(Mesa, Mesa+utility library, HEAS). For RQ5, we verify that the same
\texttt{metrics\_episode()} signature and call-site framework code remain
unchanged when switching from ecological to enterprise models, or when extending
from small-scale to large-scale evaluation.

================================================================================
EDIT 4: Add 100% agreement caveat to §5.3 (Tournament Validation)
================================================================================
LOCATION: After line 519 (after "...all tested episode budgets (4, 10, 25, 50, 100 per scenario);")
ACTION: insert
TEXT:

The 100\% voting-rule agreement likely reflects clear dominance on these
low-dimensional problems: with only 2--4 policy genes and 8 comparison scenarios,
the Pareto landscape offers unambiguous winners. A near-tie sensitivity study
(S9) confirms that this unanimity disappears when candidate policies have
statistically indistinguishable means.

================================================================================
EDIT 5: Add HV reference point to §5.2 (Multi-Scale Reproducibility)
================================================================================
LOCATION: Line 478-480 (paragraph beginning "At small scale...")
ACTION: replace
OLD:
At small scale
($n\!=\!30$, pop\,=\,20, ngen\,=\,10, steps\,=\,150),
ecological HV\,=\,7.665\,(std\,3.518, CI\,[6.424,\,8.914]).

NEW:
At small scale
($n\!=\!30$, pop\,=\,20, ngen\,=\,10, steps\,=\,150),
ecological HV\,=\,7.665\,(std\,3.518, CI\,[6.424,\,8.914]) (reference point:
mean\,biomass\,=\,-150, CV\,=\,0.5).

================================================================================
EDIT 6: Add HV reference point to §5.4 (Algorithm Agnosticism)
================================================================================
LOCATION: Line 570-572 (paragraph beginning "Scale sensitivity is reported...")
ACTION: replace
OLD:
Scale sensitivity is reported in Figure~\ref{fig:heatmap}, which maps mean HV
across a 3\,$\times$\,3 grid of steps\,$\in$\,\{140, 300, 500\} and episodes per
evaluation\,$\in$\,\{5, 10, 20\} ($n\!=\!10$ runs each).

NEW:
Scale sensitivity is reported in Figure~\ref{fig:heatmap}, which maps mean HV
across a 3\,$\times$\,3 grid of steps\,$\in$\,\{140, 300, 500\} and episodes per
evaluation\,$\in$\,\{5, 10, 20\} ($n\!=\!10$ runs each; reference point:
mean\,biomass\,=\,-200, CV\,=\,0.5).

================================================================================
EDIT 7: Add scope limitation to §5.5 (Cross-Domain Portability)
================================================================================
LOCATION: After line 628 (after "...framework modification.")
ACTION: insert
TEXT:

The portability validation is scoped to coupled ODE/tabular dynamics: both the
ecological Arena and enterprise Arena use closed-form or simple numerical
integration. Independent evaluation in spatially-explicit agent-based models,
network-based systems, or heterogeneous-agent interaction patterns remains
future work and lies outside HEAS's current design scope.

================================================================================
EDIT 8: Add artifact access statement to Conclusion
================================================================================
LOCATION: Lines 696-698 (paragraph beginning "For double-blind review...")
ACTION: replace
OLD:
For double-blind review, repository access details are omitted. An anonymized
artifact package containing all scripts under \texttt{experiments/} is prepared
for reviewers.

NEW:
For double-blind review, the repository URL is omitted from this version. Upon
acceptance, the full codebase including all experiments will be made available
at: [repository placeholder --- for camera-ready: GitHub URL]. An anonymized
artifact package containing all scripts under \texttt{experiments/} and raw
result data is prepared for reviewers.

================================================================================
SUMMARY OF CHANGES
================================================================================

Total insertions: 8 locations
Total lines added: ~170 words (~425 additional lines when formatted)
Total deletions: 0
Total replacements: 2

Critical items addressed:
  ✓ ITEM 1: Scope reframing (EDIT 7 — scope limitation statement)
  ✓ ITEM 2: Artifact access (EDIT 8 — GitHub placeholder)
  ✓ ITEM 3: Design constraints subsection (EDIT 2 — full §3.5)
  ✓ ITEM 4: NSGA-II justification (EDIT 2 — part of §3.5)
  ✓ ITEM 5: RQ1/RQ5 concreteness (EDIT 3 — falsification criteria)
  ✓ ITEM 6: HV reference points (EDITS 5–6 — explicit points in RQ2 & RQ4)
  ✓ ITEM 7: 100% agreement caveat (EDIT 4 — acknowledgment of low dimensionality)
  ✓ ITEM 8: Silent divergence example (EDIT 1 — concrete biomass scenario)
  ✓ ITEM 9: Noise stability synthetic proxy (already explicit in text, line 523)
  ✓ ITEM 10: n=20 vs n=30 justification (already justified, reference line 487)

All edits preserve existing narrative arc and maintain paper's structural integrity.

================================================================================
